\chapter{Definitions}
%\addcontentsline{toc}{chapter}{Definitions}

Following terms are used in this work:\par\bigskip

\begin{description}[style=multiline, labelwidth=5cm, leftmargin=5.5cm]

  \item[Artificial Intelligence\\in Education] 
        The interdisciplinary application of AI algorithms, machine learning models, and automated reasoning techniques to augment pedagogical methodologies and optimize individualized learning trajectories.
    
  \item[Automated Question\\Generation] 
        The computational synthesis of valid assessment items — including stems and pedagogically plausible distractors — derived directly from an underlying, unstructured knowledge base.
    
  \item[Large Language\\Model] 
        A high-parameter artificial neural network, typically leveraging a transformer topology, pre-trained on expansive corpora to computationally model syntax, semantics, and generative dialogue across natural languages.
    
  \item[Multiple-Choice\\Question] 
        An objective evaluation metric necessitating the respondent to identify the singular correct proposition against concurrently presented erroneous alternatives.
    
  \item[Natural Language\\Processing] 
        The subfield of artificial intelligence dedicated to bridging the semantic disparity between human linguistic expression and computational data interpretation, thereby enabling parsing, generation, and semantic derivation.
    
  \item[Spaced Repetition\\System] 
        An algorithmic, cognitive-science-based scheduling framework that geometrically distributes learning review intervals to disrupt the Ebbinghaus forgetting curve, subsequently maximizing long-term memory consolidation.

\end{description}